\expandafter\gdef\csname SNside1\endcsname{\node[anchor=north west,inner sep=0pt,text width=48mm,font=\SideFont,align=left] at ([xshift=138mm,yshift=-513.45pt]current page.north west) {\textcolor{SNink}{\textbf{最小反例 · 串行不等于慢}}\par\vspace{1.5mm}一条高速串行链路与一条低频8位并行总线，谁带宽更大还需各自的传送率。位宽只决定每次能送多少，线路频率、编码和协议开销同样影响结果，不能用串/并行名称代替计算。};}
\expandafter\gdef\csname SNside2\endcsname{\node[anchor=north west,inner sep=0pt,text width=48mm,font=\SideFont,align=left] at ([xshift=138mm,yshift=-364.09pt]current page.north west) {\textcolor{SNink}{\textbf{课程连接 · 把公共通路看成共享资源}}\par\vspace{1.5mm}操作系统的互斥也处理多个主体竞争一个资源。总线仲裁对应使用权分配，定时协议对应获得资源后的交互；即使仲裁公平，慢从设备仍可使某次事务长时间等待。};}
\expandafter\gdef\csname SNside3\endcsname{\node[anchor=north west,inner sep=0pt,text width=48mm,font=\SideFont,align=left] at ([xshift=138mm,yshift=-95.00pt]current page.north west) {\textcolor{SNink}{\textbf{最小反例 · 峰值相同，事务速率可不同}}\par\vspace{1.5mm}两条总线数据拍都每ns传8 B，一条每块额外等待100 ns，另一条只等1 ns；二者峰值相同，短块的有效速率却相差很大。分母是否包含等待必须先由题问确定。\par\smallskip\begin{tikzpicture}[x=1mm,y=1mm,>=Stealth,node distance=5mm]
\node[draw=SNink!50,fill=SNblue,rounded corners=1mm,text width=34mm,align=center,inner sep=1.2mm,font=\fontsize{7.5}{9}\selectfont] (m0) at (0,0) {峰值8B/ns};
\node[draw=SNink!50,fill=SNblue,rounded corners=1mm,text width=34mm,align=center,inner sep=1.2mm,font=\fontsize{7.5}{9}\selectfont] (m1) at (0,-11) {等待100ns};
\node[draw=SNink!50,fill=SNblue,rounded corners=1mm,text width=34mm,align=center,inner sep=1.2mm,font=\fontsize{7.5}{9}\selectfont] (m2) at (0,-22) {短块变慢};
\draw[->,SNink] (m0) edge[bend left=40] (m2);
\draw[->,SNink] (m1) -- (m2);
\end{tikzpicture}
\par{\fontsize{7.5}{9}\selectfont 相同数据拍峰值，等待改变事务有效速率}};}
\expandafter\gdef\csname SNside4\endcsname{\node[anchor=north west,inner sep=0pt,text width=48mm,font=\SideFont,align=left] at ([xshift=138mm,yshift=-95.00pt]current page.north west) {\textcolor{SNink}{\textbf{不变量 · 先数有效数据拍}}\par\vspace{1.5mm}32 B块经64位数据线传送，需要4个有效数据单位。突发减少的是地址、准备等重复代价，不能把32 B整块误作一个8 B数据拍；每拍几次传送还要另核时钟口径。};}
\expandafter\gdef\csname SNside5\endcsname{\node[anchor=north west,inner sep=0pt,text width=48mm,font=\SideFont,align=left] at ([xshift=138mm,yshift=-95.00pt]current page.north west) {\textcolor{SNink}{\textbf{错项排除 · 同步允许等待状态}}\par\vspace{1.5mm}共同的采样时钟只规定状态在何时被采样，不保证每个设备同样快。若协议支持等待状态，从设备可延长一次操作；``同步所以每事务固定一周期''把两项约定错误合并了。};}
\expandafter\gdef\csname SNside6\endcsname{\node[anchor=north west,inner sep=0pt,text width=48mm,font=\SideFont,align=left] at ([xshift=138mm,yshift=-95.00pt]current page.north west) {\textcolor{SNink}{\textbf{不变量 · 撤销也要确认}}\par\vspace{1.5mm}全互锁中，从设备看到请求撤销后才撤应答。少看最后一步会把半互锁误作全互锁；画四个边沿并逐项标出依赖，比按波形长度猜类别可靠。};}
\expandafter\gdef\csname SNside7\endcsname{\node[anchor=north west,inner sep=0pt,text width=48mm,font=\SideFont,align=left] at ([xshift=138mm,yshift=-95.00pt]current page.north west) {\textcolor{SNink}{\textbf{错项排除 · 主设备角色随事务}}\par\vspace{1.5mm}DMA可以发起一次主存写，CPU也可能暂不占总线。主从描述谁发起并控制本次交换，并不表示某设备永远只能发送或永远只能接收。};}
\expandafter\gdef\csname SNside8\endcsname{\node[anchor=north west,inner sep=0pt,text width=48mm,font=\SideFont,align=left] at ([xshift=138mm,yshift=-95.00pt]current page.north west) {\textcolor{SNink}{\textbf{课程连接 · 公平性与平均等待}}\par\vspace{1.5mm}固定链式优先级可能让低优先级请求久等，类似固定优先级调度的饥饿。轮转起点或老化可改善公平性，但会改变控制逻辑和响应上界，不能只比较连线数。};}
\AddToShipoutPictureFG{\begin{tikzpicture}[remember picture,overlay]\ifcsname SNside\thepage\endcsname\csname SNside\thepage\endcsname\fi\end{tikzpicture}}